\documentclass{article}
\usepackage{import}


\title{Advanced Algorithms\\Homework 06}
\author{Felix Céard-Falkenberg\\7174020}

\usepackage{../homework} % See homework.sty %

\begin{document}

\section{Question 5}

\subsection{\textbf{Prove that $C$ is a vertex cover of $G$.}}

$C$ is the set of all nodes with at least one children in a depth-first search tree of $G$.

First, if a node $v\in V$ has degree 1, i.e. there exists only one edge $(u, v) \in A$, then every depth-first search tree will contain $v$ in the leaf and $u$ in the first layer above the leaf. Therefore, $u$ will be added to $C$.

Because the depth-first search tree will have at least one leaf for each node $v\in V$ with degree $d(v) > 1$, we know that by definition, each of these nodes will be added to $C$. Therefore, $C$ is the set of all nodes with degree at least 2, with the exception of the graph containing only two nodes, in which case $C$ contains one of the two nodes. 

Because every node has an edge to a node in $C$, we know that $C$ is a vertex cover of $G$ and therefore the algorithm is correct.

\subsection{\textbf{Show that there is a matching in $T$ with at least $r/2$ edges.}}

Let us prove that there exists a matching in $T$ with at least $r/2$ edges.

Because $T$ is a tree, we can alternate between edges that are included and excluded from the matching. Because we have $r$ non-leaf nodes in $T$, we can create a matching with at least $\lfloor r/2 \rfloor$ edges. However, if $r$ is odd, we can add one more edge to the matching, namely the edge between the leaf and its parent for each branch in the tree, and therefore create a matching with at least $r/2$ edges. If $r$ is even, then we already have a matching with at least $r/2$ edges.

\subsection{\textbf{Show that $G$ has a matching of size at least $|C|/2$}.}

Let us prove that $G$ has a matching of size at least $|C|/2$.

Because $C$ is a valid vertex cover of $G$, we can create a matching of size at least $|C|/2$ by selecting one edge from each node in $C$. Because each edge is incident to at least one node in $C$, we know that the matching is valid. 

By adding the edges to the leafs --- if $|C|$ is odd, otherwise $|C|/2$ is already guaranteed --- we can guarantee that the matching is of size at least $|C|/2$.

% Because $C$ is a vertex cover of $G$, we know that every edge in $G$ is incident to at least one node in $C$. Therefore, we can create a matching of size at least $|C|/2$ by selecting one edge from each node in $C$.

\subsection{\textbf{Prove that the algorithm outputs a vertex cover whose size is at most two times the size of a minimum vertex cover.}}

% In the best case scenario, $G$ has a minimal vertex cover of size $|C| = |V|/2$. However, in the worst case scenario, the algorithm will output $C = V$ and therefore 

This is not true I think. Consider the graph in Figure~\ref{fig:example-graph}. Either you do not include all nodes in the leafs, however, the only node not in the leaf is the root node, which, if $C = \set{D}$, is not a valid vertex cover. Otherwise, in this case, the algorithm would return $C = V$ although the minimal vertex cover is $C = \set{E}$. Because $|C| = 1 < 2 \cdot |C| < |V|$, I think this is not true, or the definition is ambiguous.

\begin{figure}[H]
  \centering
  \includegraphics[width=0.5\linewidth]{graph.png}
  \caption{Left: An example graph. Right: an excerpt of the depth-first search tree.}
  \label{fig:example-graph}
\end{figure}


\subsection{\textbf{The algorithm above was stated for connected graphs. What would you do for graphs that are not connected?}}

Start the search for each node not already found. During the DFS, you keep a set of all checked nodes $K$ and if $|K| < |V|$, you start a new DFS from a node $v\in V \setminus K$.





\end{document}
